\documentclass[11pt, a4paper]{article}

\usepackage[T1]{fontenc}
\usepackage[utf8]{inputenc}
\usepackage[french]{babel}
\usepackage[margin=1.4cm,top=1.3cm,bottom=1.3cm]{geometry}
\usepackage{amsmath, amssymb}
\usepackage{booktabs}
\usepackage{enumitem}
\usepackage{titlesec}
\usepackage[expansion=false]{microtype}
\usepackage{xcolor}
\usepackage{graphicx}
\usepackage{tikz}
\usetikzlibrary{arrows.meta, positioning, shapes.geometric}
\usepackage[hidelinks]{hyperref}
\usepackage{multicol}
\usepackage{float}

\setlength{\parskip}{0.2em}
\setlength{\parindent}{0pt}
\titlespacing*{\section}{0pt}{0.5em}{0.2em}
\setlist{nosep, leftmargin=1.5em}
\renewcommand{\baselinestretch}{0.97}

\title{\vspace{-1.5em}\textbf{Analyse exp\'erimentale : IA Responsable\\sur le German Credit Dataset}}
\author{IA708 -- T\'el\'ecom Paris, 2026}
\date{}

\begin{document}
\maketitle
\vspace{-2em}

\section{Protocole exp\'erimental}

Le jeu UCI German Credit (Hofmann, 1994) contient 1\,000 demandes de cr\'edit, 20 attributs (6 num\'eriques, 14 cat\'egoriels), cible binaire d\'efaut/non-d\'efaut (30\,\% / 70\,\%). Deux attributs sensibles : \textbf{genre} (extrait de \texttt{personal\_status\_sex}, 690 hommes / 310 femmes) et \textbf{\^age} binaris\'e au seuil 25 ans (810 older / 190 younger). Les attributs sensibles sont initialement retir\'es des features ; nous discutons ce choix \`a la lumi\`ere d'une analyse de proxies (\S 4).

\paragraph{Mod\`ele.} R\'egression logistique L2 (\texttt{sklearn.LogisticRegression}, \texttt{liblinear}, $C=1$, \texttt{max\_iter}\,=\,1000). \textbf{Convergence} v\'erifi\'ee : log-loss train atteint 0.43 d\`es 10 it\'erations, val 0.56 (proche de la r\'ef\'erence 0.5 attendue sur ce d\'es\'equilibre).

\paragraph{Donn\'ees.} Split stratifi\'e 60/20/20 (train/val/test, \texttt{random\_state=42}). Pr\'etraitement \texttt{sklearn.ColumnTransformer} : \texttt{StandardScaler} pour les num\'eriques, \texttt{OneHotEncoder} pour les cat\'egoriels (56 colonnes finales). Le seuil de classification est calibr\'e sur la val pour maximiser la balanced accuracy ($\tau = 0.18$).

\paragraph{Outils.} \texttt{scikit-learn} (mod\`ele, m\'etriques, calibration, permutation importance), \texttt{fairlearn} (\texttt{MetricFrame}, \texttt{demographic\_parity\_difference}, \texttt{equalized\_odds\_difference}), \texttt{shap} (\texttt{LinearExplainer}). Reproductibilit\'e par \texttt{papermill}.

\section{M\'ethodes d'\'equit\'e}

\paragraph{Pr\'e-traitement (reweighing).} Poids $w_i = P(S{=}s_i)\,P(Y{=}y_i)\,/\,P(S{=}s_i, Y{=}y_i)$ rendant $S \perp Y$ dans la distribution pond\'er\'ee.

\paragraph{Post-traitement (seuils par groupe).} On cherche un seuil $\tau_g$ par groupe sensible \'egalisant le taux de s\'election cible (Demographic Parity), calibr\'e sur la validation puis appliqu\'e au test.

\section{R\'esultats principaux}

\begin{table}[ht]
\centering
\small
\begin{tabular}{llcccc}
\toprule
\textbf{Attribut} & \textbf{Mod\`ele} & \textbf{AUC} & \textbf{BalAcc} & $|\Delta\text{DP}|$ & $|\Delta\text{EO}|$ \\
\midrule
Genre & Baseline        & 0.785 & 0.669 & 0.010 & 0.058 \\
Genre & Reweighing      & 0.785 & 0.665 & 0.006 & 0.050 \\
Genre & Baseline + PP   & 0.785 & 0.669 & 0.105 & 0.196 \\
Genre & Reweighing + PP & 0.785 & 0.669 & 0.128 & 0.230 \\
\midrule
\^Age  & Baseline        & 0.785 & 0.669 & 0.125 & 0.066 \\
\^Age  & Reweighing      & 0.784 & 0.640 & 0.132 & 0.098 \\
\^Age  & Baseline + PP   & 0.785 & 0.644 & 0.060 & 0.135 \\
\^Age  & Reweighing + PP & 0.784 & 0.650 & 0.069 & 0.221 \\
\bottomrule
\end{tabular}
\caption{Performance et \'equit\'e sur le test (PP = post-processing par seuils par groupe).}
\label{tab:results}
\end{table}

\paragraph{Genre.} Apr\`es retrait de \texttt{personal\_status\_sex}, le biais direct dispara\^it ($|\Delta\text{DP}| = 0.010$). Le reweighing n'a pas d'effet, mais le post-processing \textbf{cr\'ee} un biais artificiel ($|\Delta\text{DP}| = 0.105$, $|\Delta\text{EO}| = 0.196$) : sur un d\'es\'equilibre 690/310 et un test de 200, calibrer des seuils par groupe sur 200 exemples de validation se g\'en\'eralise mal.

\paragraph{\^Age.} Le biais reste significatif apr\`es retrait ($|\Delta\text{DP}| = 0.125$), coh\'erent avec les taux de d\'efaut tr\`es diff\'erents (older 27\,\%, younger 42\,\%). Le post-processing r\'eussit \`a r\'eduire $|\Delta\text{DP}|$ \`a 0.060 mais \textbf{double} $|\Delta\text{EO}|$ (0.066 vers 0.135), illustration empirique du \textbf{th\'eor\`eme d'impossibilit\'e} de Chouldechova~\cite{chouldechova2017} et Kleinberg et al.~\cite{kleinberg2016}.

\section{Interpr\'etabilit\'e et d\'etection de proxies}

Pour la r\'egression logistique, $\phi_i(x) = w_i (x_i - \mathbb{E}[x_i])$ est exact (Lundberg \& Lee~\cite{lundberg2017}, Corollary~1). SHAP et permutation importance concordent sur le top : \texttt{checking\_status}, \texttt{credit\_history}, \texttt{credit\_amount}, \texttt{duration\_in\_month}.

\paragraph{Proxies multivari\'es.} Suivant la recommandation de C.~Laclau, on entra\^ine une r\'egression logistique \texttt{features $\to$ attribut sensible} : \textbf{AUC = 0.86 pour l'\^age}, AUC = 0.69 pour le genre. L'attribut sensible est largement reconstructible depuis les autres features. \textbf{Cons\'equence} : retirer l'attribut sensible des inputs ne suffit pas \`a garantir l'\'equit\'e ; les m\'ethodes "fairness through awareness" (le conserver pour mitiger explicitement) seraient th\'eoriquement plus efficaces.

\paragraph{SHAP avant vs apr\`es retrait.} Mod\`ele \textbf{avec} attributs sensibles : AUC=0.781 ; \textbf{sans} (baseline) : AUC=0.785 ($\Delta$ n\'egligeable). $\Delta$ Mean |SHAP| (sans $-$ avec) montre que les features absorbent l'effet : \texttt{present\_employment\_since} (+0.036), \texttt{telephone} (+0.022), \texttt{credit\_history} (+0.020). Confirmation graphique de l'AUC proxy=0.86 (figure dans le notebook et le beamer).

\section{Quantification d'incertitude (substitut \`a la robustesse adversariale)}

Sur conseil de l'enseignante, la robustesse a \'et\'e remplac\'ee par une \textbf{quantification d'incertitude par bootstrap} (cf. cours \textit{Robustness and Uncertainty}). On entra\^ine $B = 200$ mod\`eles sur des \'echantillons tir\'es avec remplacement, et on calcule des intervalles de confiance bootstrap sur toutes les m\'etriques :

\begin{table}[ht]
\centering
\small
\begin{tabular}{lcccc}
\toprule
\textbf{M\'etrique} & \textbf{Moyenne} & \textbf{\'Ecart-type} & \textbf{IC 2.5\,\%} & \textbf{IC 97.5\,\%} \\
\midrule
AUC          & 0.764 & 0.019 & 0.723 & 0.795 \\
BalAcc       & 0.682 & 0.023 & 0.637 & 0.726 \\
$|\Delta\text{DP}|$ genre & 0.035 & 0.028 & 0.001 & 0.105 \\
$|\Delta\text{EO}|$ genre & 0.089 & 0.044 & 0.017 & 0.176 \\
$|\Delta\text{DP}|$ \^age  & 0.104 & 0.046 & 0.020 & 0.202 \\
$|\Delta\text{EO}|$ \^age  & 0.086 & 0.048 & 0.018 & 0.194 \\
\bottomrule
\end{tabular}
\caption{Intervalles de confiance bootstrap (B=200, IC 95\%).}
\end{table}

\textbf{L'IC95\,\% de $|\Delta\text{DP}|$ \^age s'\'etend de 0.02 \`a 0.20}. Le biais empirique de 0.125 est dans cet intervalle mais l'incertitude est tr\`es large (test de 200), ce qui plaide pour la \textbf{prudence m\'ethodologique} : sur ce dataset, les \'ecarts d'\'equit\'e ne sont pas significatifs \`a un niveau pr\'ecis.

\section*{Conclusion et limites}

L'effet des corrections d'\'equit\'e d\'epend fortement de la structure des proxies : reweighing inefficace pour le genre (proxies faibles, AUC proxy = 0.69), modificateur pour l'\^age (proxies forts, AUC proxy = 0.86). Le th\'eor\`eme d'impossibilit\'e est confirm\'e empiriquement. Le bootstrap montre que sur un dataset de 1\,000 lignes, les conclusions sur l'\'equit\'e doivent \^etre nuanc\'ees par l'incertitude statistique (IC dépassant 0.15).

\textbf{Limites} : approche lin\'eaire, reweighing simple, pas de robustesse adversariale (substitution par bootstrap). \textbf{Pistes} : essayer \texttt{fairlearn.ExponentiatedGradient} ou un transport optimal \`a la place du reweighing ; conserver l'\^age et appliquer une mitigation in-processing.

\textbf{Outils d'assistance.} Des LLM (Claude, GPT) ont \'et\'e utilis\'es pour l'aide \`a la r\'edaction du code (pipeline \texttt{sklearn}, int\'egration \texttt{fairlearn}/\texttt{shap}) et la structuration. Toutes les d\'ecisions m\'ethodologiques (choix des m\'etriques, interpr\'etation, conception du bootstrap) ont \'et\'e prises et discut\'ees en \'equipe.

\paragraph{\small Références} {\scriptsize
A.~Chouldechova, \textit{Fair prediction with disparate impact}, Big Data 5(2), 2017, \href{https://arxiv.org/abs/1703.00056}{arXiv:1703.00056}\,;
J.~Kleinberg et al., \textit{Inherent trade-offs in the fair determination of risk scores}, ITCS 2017, \href{https://arxiv.org/abs/1609.05807}{arXiv:1609.05807}\,;
S.~Lundberg, S.-I.~Lee, \textit{A Unified Approach to Interpreting Model Predictions}, NeurIPS 2017, \href{https://arxiv.org/abs/1705.07874}{arXiv:1705.07874} (Linear SHAP, Corollary~1, p.~6)\,;
F.~Kamiran, T.~Calders, \textit{Data preprocessing techniques for classification without discrimination}, KAIS 33(1), 2012, \href{https://doi.org/10.1007/s10115-011-0463-8}{doi:10.1007/s10115-011-0463-8}\,;
R.~K.~E.~Bellamy et al., \textit{AI Fairness 360}, IBM 2018, \href{https://arxiv.org/abs/1810.01943}{arXiv:1810.01943}\,;
M.~Hardt, E.~Price, N.~Srebro, \textit{Equality of Opportunity in Supervised Learning}, NeurIPS 2016, \href{https://arxiv.org/abs/1610.02413}{arXiv:1610.02413}\,;
B.~Efron, \textit{Bootstrap methods}, Annals of Statistics 7(1), 1979, \href{https://doi.org/10.1214/aos/1176344552}{doi:10.1214/aos/1176344552}\,;
C.~Guo et al., \textit{On Calibration of Modern Neural Networks}, ICML 2017, \href{https://arxiv.org/abs/1706.04599}{arXiv:1706.04599}\,;
S.~Bird et al., \textit{Fairlearn}, Microsoft 2020, \href{https://fairlearn.org/}{fairlearn.org}.
}

\end{document}
